\section{Conclusion}
\label{sec:conclusion}

We presented \savyar{}, a Turkish morphological disambiguator that ranks the output of a rule-based analyzer without seeing open-class roots. Root-specific rules are applied when candidates are generated, and the ranker chooses the suffix sequence of each word in the context of all analyses of the sentence. Because the ranker never sees the root, unambiguous training words can be presented with the analyses of ambiguous words, which adds \TBD{synth-gain} points on ambiguous words. On TrMor2018, under the protocol of \citet{akyurek-etal-2019-morphological}, \savyar{} reaches \TBD{acc-all}\% on all words and \TBD{acc-amb}\% on ambiguous words. The analyzer proposes the gold analysis for \TBD{coverage}\% of words, so \TBD{coverage-share}\% of the remaining errors lie in the analyzer rather than the ranker. Sentence context adds \TBD{context-gain} points on ambiguous words, and giving the ranker access to roots changes accuracy by \TBD{root-delta} points.

Two directions follow from these results. Mapping \savyar{} analyses back to Universal Dependencies features would allow evaluation on the UD Turkish treebanks used in the SIGMORPHON 2019 shared task \citep{mccarthy-etal-2019-sigmorphon} and by \citet{can-etal-2022-joint}. Splitting the root token by finer categories, for example separating adjectives and proper nouns from other nominals, would address the distinctions that the current encoding cannot express. \TBD{revise once the root-aware result is known}.
